% ---------------------------------------------------------------------
%  08 · Exigences
%
%    \begin{exigence}[type=EF, priorite=haute]{Titre}
%      Description…
%      \label{ex:scan}
%    \end{exigence}
%
%  Identifiant auto : <type>-<chapitre>.<n>  (ex. EF-03.02), \ref-able.
%  type     : EF (fonctionnelle) · ENF (non fonctionnelle) · libre
%  priorite : haute · moyenne · basse
% ---------------------------------------------------------------------

\newcounter{exigence}[chapter]

\makeatletter
\pgfkeys{
  /exigence/.is family, /exigence,
  type/.store in     = \exg@type,
  priorite/.store in = \exg@prio,
  default/.style     = {type=EF, priorite=haute},
}

\newcommand\exg@badge[1]{%
  \ifdefstring{#1}{haute}  {\badge[herbe]{Priorité haute}}{%
  \ifdefstring{#1}{moyenne}{\badge[agrume]{Priorité moyenne}}{%
  \ifdefstring{#1}{basse}  {\badge[contour]{Priorité basse}}{%
                            \badge[sauge]{#1}}}}%
}

\renewcommand\theexigence{\exg@type-\deuxchiffres{\value{chapter}}.\deuxchiffres{\value{exigence}}}

% Un compteur par type (créé au premier usage) : EF-01.01, ENF-01.01, …
\newcommand\exg@compter{%
  \ifcsname c@exg@\exg@type\endcsname\else
    \newcounter{exg@\exg@type}[chapter]%
  \fi
  \stepcounter{exg@\exg@type}%
  \setcounter{exigence}{\numexpr\value{exg@\exg@type}-1\relax}%
  \refstepcounter{exigence}%
}

\NewDocumentEnvironment{exigence}{O{} m}{%
  \pgfkeys{/exigence, default, #1}%
  \exg@compter
  \begin{tcolorbox}[marche]
    \noindent
    \textlabel[action]{\theexigence}\hfill\exg@badge{\exg@prio}\par
    \vspace{0.4em}%
    {\policeIntertitre\large #2\par}%
    \vspace{0.2em}%
}{%
  \end{tcolorbox}%
}
\makeatother
